\slajdid{02-s035}
\begin{frame}[label=02-s035]
\gusttekst
\begin{columns}[c,onlytextwidth]
\begin{column}{.27\textwidth}
Tabela sa 3 promenljive
\end{column}
\begin{column}{.29\textwidth}\centering
\begingroup
\fontsize{14.52}{16.13}\selectfont
\renewcommand{\small}{\fontsize{14.52}{16.13}\selectfont}
\scalebox{.62}{%
\begin{karnaugh-map}(label=middle)[2][4][1][$A$][$B$][$C$]
\foreach \x/\y/\lab in {1/3/0,2/3/1,1/2/2,2/2/3,1/1/6,2/1/7,1/0/4,2/0/5}{
 \node[anchor=south east,inner sep=2pt,font=\fontsize{12.91}{14}\selectfont] at (\x,\y) {\lab};
}
\end{karnaugh-map}}
\endgroup

\end{column}
\begin{column}{.05\textwidth}\centering ili\end{column}
\begin{column}{.35\textwidth}\centering
\begingroup
\fontsize{14.52}{16.13}\selectfont
\renewcommand{\small}{\fontsize{14.52}{16.13}\selectfont}
\scalebox{.62}{%
\begin{karnaugh-map}(label=middle)[4][2][1][$A$][$B$][$C$]
\foreach \x/\y/\lab in {1/1/0,2/1/1,3/1/3,4/1/2,1/0/4,2/0/5,3/0/7,4/0/6}{
 \node[anchor=south east,inner sep=2pt,font=\fontsize{12.91}{14}\selectfont] at (\x,\y) {\lab};
}
\end{karnaugh-map}}
\endgroup

\end{column}
\end{columns}

Treba uočiti da su u 1. slučaju po definiciji susedna polja sa indeksima 0 i 4, odnosno 1 i 5, što vizuelno odgovara presavijanju tabele oko horizontalne ose po srediti i pravljenjem prstena.

U 2. Slučaju su susedna i polja 0 i 2, odnosno 4 i 6 što se opet vizuelno dobija presavijanjem tabele oko vertikalne ose i pravljenjem prstena.
\end{frame}
